\largerpage
\sscp{Labial plosives}{06-01-04}
Additional evidence for Proto-Aru comes from the sound changes
affecting the labial plosives.
We begin with the changes affecting PMP *p.
The changes affecting PMP *p are simplest in Ujir,
Kompane, \tsc{Kola}, and Lola.
Ujir has PMP *p > {\0} initially
and PMP *p > \it{w,} {\0} medially.
\tsc{Kola} has PMP *p > \it{f} in all positions.
Kompane has PMP *p > \it{h} initially,
and PMP *p > \it{h,} {\0} word medially.
Lola has initial PMP *p > \it{h},
but medial PMP *p > {\0},
\it{w, h,} or rounding of an adjacent vowel \tsc{(+round)}.
The changes affecting PMP *p in the \tsc{South Aru} languages
are more complex, with several splits attested.
The different reflexes are summarised in \trf{06-10}.
They include: PMP *p > \it{w},
fortition (of intermediate **w),
loss, and rounding of an adjacent vowel.
\citet{Nivens2017} posits Proto-South Aru **β to account for these reflexes.
Examples of PMP *p in Aru are given in \trf{06-11}.

\begin{table}
	\centering\tcp{Reflexes of PMP *p in \tsc{South Aru}}{06-10}
	\st{6pt}\begin{tabular}{llllll}\lsptoprule
Lorang	&		&	w	&	gw	&		&	{\0}\\
Dobel	&		&	w	&	kw	&		&	{\0}\\
Koba	&		&	o	&	ʍ	&		&	{\0}\\
Barakai	&		&		&		&		&	{\0}\\
Karey	&		&	β	&	g	&	\tsc{+round}	&	{\0}\\
Batuley	&		&	w	&		&		&	{\0}\\
Mariri	&		&	w	&		&		&	{\0}\\
Manombai	&		&		&		&		&	{\0}\\
\tsc{Tarangan}	&		&		&		&	\tsc{+round}	&	{\0}\\
		\lspbottomrule
	\end{tabular}
\end{table}

\begin{table}\tcp{\tsc{Aru} *p > **ɸ}{06-11}
\begin{adjustbox}{width=\textwidth}
	\begin{threeparttable}\st{2pt}
	\begin{tabular}{llllllll}\lsptoprule
local gl.	&	{\hr}fire	&	{\hr}bat	&	{\hr}thin	&	{\hr}hungry	&	{\hr}tooth	&	{\hr}shoot	&	{\hr}ten\\
PMP	&	\hr*ha\tbr{p}uy	&	\hr*\tbr{p}aniki	&	\hr*ma-ni\tbr{p}is	&	\hr*la\tbr{p}aR	&	\hr*[nŋ]i\tbr{p}ən	&	\hr*\tbr{p}anaq	&	\\
Proto-Aru	&	**a\tbr{ɸ}uy	&	**\tbr{ɸ}aniki	&	**mani\tbr{ɸ}i-	&	**kan-la\tbr{ɸ}aR	&	**[nŋ]e\tbr{ɸ}in	&	**-\tbr{ɸ}ana	&	**\tbr{ɸ}uR\su{Δ}\\ \midrule
Ujir	&	{\hr\hr}ʤe	&		&	{\hr\hr}maníʔ	&	{\hr\hr}kana\tbr{w}ay	&	{\hr\hr}nein	&	{\hr\hr}-an	&	{\hr\hr}uy sia\\
Kompane	&	{\hr\hr}a\tbr{h}i	&	{\hr\hr}\tbr{h}anisi	&	{\hr\hr}maniay	&	{\hr\hr}kanaa	&	{\hr\hr}ne\tbr{h}in	&		&	{\hr\hr}\tbr{h}uɸaɸi\\
East Kola	&	{\hr\hr}e\tbr{f}	&	{\hr\hr}\tbr{f}anís	&	{\hr\hr}mani\tbr{f}i	&	{\hr\hr}kanaah-ni	&	{\hr\hr}no\tbr{f}in	&	{\hr\hr}-\tbr{f}an	&	{\hr\hr}\tbr{f}uh\\
Lola	&	{\hr\hr}ʤa\tbr{h}i	&	{\hr\hr}\tbr{h}anisi	&	{\hr\hr}manuay	&	{\hr\hr}ʔaná\tbr{w}i	&	{\hr\hr}ŋein	&		&	{\hr\hr}\tbr{h}aɸi ye\\  \midrule
Lorang	&	{\hr\hr}ʤay	&	{\hr\hr}\tbr{w}eníu	&	{\hr\hr}mani\tbr{gw}ay	&		&	{\hr\hr}yɪŋɪn	&		&	{\hr\hr}uru\\
Dobel	&	{\hr\hr}say	&	{\hr\hr}\tbr{w}oyníw	&	{\hr\hr}mani\tbr{kw}ay	&	{\hr\hr}na-ka‹l›la\tbr{w}ar-ni	&	{\hr\hr}ŋeyan-	&	{\hr\hr}-\tbr{w}ana	&	{\hr\hr}\tbr{w}ur\\
Koba	&	{\hr\hr}say	&	{\hr\hr}\tbr{o}éni	&	{\hr\hr}maní\tbr{ʍ}ey	&		&		&		&	{\hr\hr}ur\\
Barakai	&	{\hr\hr}aa	&	{\hr\hr}aaní	&	{\hr\hr}mania	&	{\hr\hr}kanáir	&	{\hr\hr}neen	&		&	{\hr\hr}urwéo\\
Karey	&	{\hr\hr}aa	&	{\hr\hr}anido	&	{\hr\hr}manua	&	{\hr\hr}kaná\tbr{β}ir	&	{\hr\hr}ne\tbr{g}en	&		&	{\hr\hr}orwéw\\
Batuley	&	{\hr\hr}ʤekáy	&	{\hr\hr}anís	&	{\hr\hr}mani\tbr{w}e	&	{\hr\hr}kana\tbr{w}er	&	{\hr\hr}neyen	&		&	{\hr\hr}ur\\
Mariri	&	{\hr\hr}ay	&	{\hr\hr}adíʧ	&	{\hr\hr}madiway	&		&	{\hr\hr}deen	&		&	{\hr\hr}orɸéɸ\\
Manombai	&	{\hr\hr}a\tbr{w}	&	{\hr\hr}nisi	&	{\hr\hr}maniay	&	{\hr\hr}kanári	&	{\hr\hr}nini	&	{\hr\hr}-en\su{‡}	&	{\hr\hr}uraɸa ia\\
NE. Tarang.	&	{\hr\hr}ɛ	&	{\hr\hr}ɔnít	&	{\hr\hr}manuɛ	&		&	{\hr\hr}nin	&		&	{\hr\hr}urɸeɸa\\
SE. Tarang.	&	{\hr\hr}ɛy	&	{\hr\hr}ɔnít	&	{\hr\hr}mɛnugway	&	{\hr\hr}kanagwir	&	{\hr\hr}nɛin	&		&	{\hr\hr}ɔrɸɛɸɔ\\
NW. Tar.	&	{\hr\hr}ɔ	&	{\hr\hr}nik	&	{\hr\hr}manuay	&	{\hr\hr}kanɔir	&	{\hr\hr}nin	&	{\hr\hr}-ɔn	&	{\hr\hr}ur\\
SW. Tar.	&	{\hr\hr}ɔ\tbr{w}	&	{\hr\hr}aník\su{†}	&	{\hr\hr}mɛnuay	&	{\hr\hr}kanɔir	&	{\hr\hr}nin	&	{\hr\hr}-ɔn	&	{\hr\hr}ur\\
		\lspbottomrule
	\end{tabular}
		\begin{tablenotes}\tnsize
			\item[†]	{Southwest Tarangan \it{aník} ‘fruit bat’ is from Feruni and Kabalukin villages.}
			\item[‡]	{Manombai \it{-en} ‘shoot’ is from Gardakau village.}
			\item[Δ]	{Several reflexes of **ɸuR ‘ten’ here are from a compound **ɸuR-papay.}
		\end{tablenotes}
	\end{threeparttable}
\end{adjustbox}
\end{table}

PMP *b > \it{ɸ} in most Aru languages.
\tsc{Tarangan} usually has *b > \it{p} when it becomes word final
due to loss of an earlier vowel,
and *b > \it{ɸ {\tl} ʍ {\tl} h} in other word positions with
the result varying by the variety of \tsc{Tarangan} and/or by
quality of the following vowels.
Karey has *b > ɸ /{\gap}o
and *b > \it{w} before other vowels.
Barakai has *b > \it{w}.

While Kola has *b > \it{ɸ} in all word positions the phoneme /ɸ/ has allophony.
\citet{Nivens2017} describes Kola /ɸ/ as “(always?) rounded,
and sometimes half-voiced”, while \citet[34]{TakataTakata1992}
describe free variation between a voiceless
and voiced bilabial fricative,
/ɸ/ → [ɸ] {\tl} [β].
We transcribe this segment ‹ɸʷ›.
Kola bilabial /ɸ/ contrasts with labio-dental /f/.

\citet{Nivens2017} reconstructs the reflex of PMP *b in Proto-Aru
as **p, noting that an alternate reconstruction would be **β.
Examples of PMP *b > **p are given in \trf{06-12}.

\begin{table}\tcp{\tsc{Aru} *b > **p}{06-12}
\begin{adjustbox}{width=\textwidth}
	\begin{threeparttable}\st{2pt}
	\begin{tabular}{llllllll}\lsptoprule
local gl.	&	{\hr}body hair	&	{\hr}moon	&	{\hr}fruit	&	{\hr}stone	&	{\hr}mouth	&	{\hr}top	&	{\hr}pig\\
PMP	&	\hr*\tbr{b}ulu	&	\hr*\tbr{b}ulan	&	\hr*\tbr{b}uaq	&	\hr*\tbr{b}atu	&	\hr*\tbr{b}aq\tbr{b}aq	&	\hr*\tbr{b}u\tbr{b}uŋ	&	\hr*\tbr{b}a\tbr{b}uy\\
Proto-Aru	&	**\tbr{p}ulu-	&	**\tbr{p}ulan	&	**\tbr{p}ua-\su{‡}	&	**\tbr{p}atu	&	**\tbr{p}a\tbr{p}a-	&	**\tbr{p}u\tbr{p}un	&	**\tbr{p}a\tbr{p}uy\\ \midrule
Ujir	&	{\hr\hr}\tbr{ɸ}ul	&	{\hr\hr}\tbr{ɸ}ulan	&	{\hr\hr}\tbr{ɸ}uay	&	{\hr\hr}\tbr{ɸ}at	&	{\hr\hr}\tbr{ɸ}e\tbr{ɸ}a	&		&	\\
Kompane	&		&	{\hr\hr}\tbr{ɸ}ulan	&	{\hr\hr}\tbr{ɸ}ui	&		&	{\hr\hr}\tbr{ɸ}a\tbr{ɸ}i	&		&	\\
East Kola	&		&	{\hr\hr}\tbr{ɸʷ}ulan	&	{\hr\hr}\tbr{ɸʷ}uui	&	{\hr\hr}\tbr{ɸʷ}at	&	{\hr\hr}\tbr{ɸʷ}a\tbr{ɸʷ}a-	&	{\hr\hr}\tbr{ɸʷ}u\tbr{ɸʷ}un\su{¶}	&	{\hr\hr}\tbr{ɸʷ}e\tbr{ɸʷ}\\
Lola	&	{\hr\hr}\tbr{ɸ}uli	&	{\hr\hr}\tbr{ɸ}ulan	&	{\hr\hr}\tbr{ɸ}ua-di	&		&	{\hr\hr}\tbr{ɸ}a\tbr{ɸ}i	&	{\hr\hr}\tbr{ɸ}u\tbr{ɸ}in	&	\\
Lorang	&	{\hr\hr}\tbr{ɸ}ul	&	{\hr\hr}\tbr{ɸ}ulan	&	{\hr\hr}\tbr{ɸ}úʤa	&	{\hr\hr}\tbr{ɸ}atu	&	{\hr\hr}\tbr{ɸ}a\tbr{ɸ}a	&		&	\\
Dobel	&	{\hr\hr}\tbr{ɸ}ulu-\su{ǁ}	&	{\hr\hr}\tbr{ɸ}ulan	&	{\hr\hr}\tbr{ɸ}usi	&	{\hr\hr}\tbr{ɸ}atu	&	{\hr\hr}\tbr{ɸ}a\tbr{ɸ}a-m	&	{\hr\hr}\tbr{ɸ}u\tbr{ɸ}un	&	{\hr\hr}\tbr{ɸ}a\tbr{ɸ}i\\
Koba	&	{\hr\hr}\tbr{ɸ}ul	&	{\hr\hr}\tbr{ɸ}ulan	&	{\hr\hr}\tbr{ɸ}úsa	&	{\hr\hr}\tbr{ɸ}atu	&	{\hr\hr}\tbr{ɸ}e\tbr{ɸ}a	&	{\hr\hr}\tbr{ɸ}i\tbr{ɸ}íŋ	&	\\
Barakai	&	{\hr\hr}\tbr{w}olʊ-n	&	{\hr\hr}\tbr{w}olan	&	{\hr\hr}\tbr{w}ii	&	{\hr\hr}\tbr{w}aʔ	&	{\hr\hr}\tbr{w}eʊ	&	{\hr\hr}\tbr{w}uun	&	\\
Karey	&	{\hr\hr}\tbr{ɸ}o\tbr{ɸ}olo-n	&	{\hr\hr}\tbr{ɸ}olan	&	{\hr\hr}\tbr{w}uaw	&		&	{\hr\hr}\tbr{w}ew	&	{\hr\hr}\tbr{w}uwun	&	\\
Batuley	&	{\hr\hr}\tbr{ɸ}oloi	&	{\hr\hr}\tbr{ɸ}ulan	&	{\hr\hr}\tbr{ɸ}ui	&	{\hr\hr}\tbr{ɸ}at\su{Δ}	&	{\hr\hr}\tbr{ɸ}a\tbr{ɸ}a-ŋ	&	{\hr\hr}\tbr{ɸ}u\tbr{ɸ}un	&	{\hr\hr}\tbr{ɸ}ae\tbr{ɸ}\\
Mariri	&	{\hr\hr}\tbr{ɸ}ul\tbr{ɸ}ulu-n	&	{\hr\hr}\tbr{ɸ}ulan	&	{\hr\hr}\tbr{ɸ}ui	&		&	{\hr\hr}\tbr{ɸ}e\tbr{ɸ}	&	{\hr\hr}\tbr{ɸ}u\tbr{ɸ}un	&	\\
Manombai	&		&	{\hr\hr}\tbr{ɸ}ulan	&	{\hr\hr}\tbr{ɸ}uay	&	{\hr\hr}\tbr{ɸ}atu	&	{\hr\hr}\tbr{ɸ}a\tbr{ɸ}a-y	&		&	{\hr\hr}\tbr{ɸ}a\tbr{ɸ}u\\
NE. Tar. (Fb)\su{†}	&		&	{\hr\hr}\tbr{h}ulan	&		&	{\hr\hr}\tbr{h}ɔt	&	{\hr\hr}\tbr{ɸ}e\tbr{ɸ}a	&		&	\\
NE. Tar. (Gm)\su{†}	&		&	{\hr\hr}\tbr{h}olan	&	{\hr\hr}\tbr{h}uʤa	&	{\hr\hr}\tbr{ʍ}at	&	{\hr\hr}\tbr{ʍ}e\tbr{ʍ}a	&	{\hr\hr}\tbr{ʍ}u\tbr{ʍ}in	&	\\
SE. Tar. (Sr)\su{†}	&		&	{\hr\hr}\tbr{ɸ}ôlan	&	{\hr\hr}\tbr{ɸ}uʤɔ	&	{\hr\hr}\tbr{ɸ}at	&	{\hr\hr}\tbr{ɸ}ɛ\tbr{ɸ}ɔ	&	{\hr\hr}\tbr{ɸ}u\tbr{ɸ}in	&	\\
NW. Tar. (Ju)\su{†}	&		&	{\hr\hr}\tbr{h}ulan	&	{\hr\hr}\tbr{h}uɛ	&	{\hr\hr}\tbr{h}ɔt	&	{\hr\hr}\tbr{h}a\tbr{h}e	&		&	{\hr\hr}\tbr{h}ɔ\tbr{p}\\
SW. Tar. (Kk)\su{†}	&	{\hr\hr}\tbr{ɸ}ul-di\su{ǁ}	&	{\hr\hr}\tbr{ɸ}ôlan	&	{\hr\hr}\tbr{ɸ}uay	&	{\hr\hr}\tbr{ɸ}ɔt	&	{\hr\hr}\tbr{ɸ}a\tbr{ɸ}a-y	&	{\hr\hr}\tbr{ɸ}u\tbr{ɸ}an\su{¶}	&	{\hr\hr}\tbr{ɸ}ɔ\tbr{p}\\
		\lspbottomrule
	\end{tabular}
		\begin{tablenotes}\tnsize
			\item[†]	{Northeast Tarangan (Fb) from Fatlabata,
								Northeast Tarangan (Gm) from Gomar Meti,
								Southeast Tarangan (Sr) from Salarem village, Southwest Tarangan
								(Kk) from Kalar-kalar, Northwest Tarangan (Ju) from Juring.}
			\item[‡]	{Reflexes of **pua- ‘fruit’ are \tsc{3sg} forms except Lola
								\it{ɸua-di} which is \tsc{3pl}.}
			\item[Δ]	{Batuley \it{ɸat} is given as ‘dead weight anchor; anchor
								consisting of a blunt heavy weight such as a stone or piece of concrete’ \citep[251]{Daigle2015}.}
			\item[{ǁ}]	{Dobel \it{ɸulu-}
								and Southwest Tarangan \it{ɸul-di} are ‘eyebrows, eye lashes’.}
			\item[¶]	{East Kola \it{ɸʷuɸʷun}
								and Southwest Tarangan \it{ɸuɸan} are ‘roof ridge’.
								Other reflexes mean ‘top, surface’.
								PMP *bubuŋ is reconstructed meaning ‘ridge of the roof; ridge
								of a mountain, peak; deck of a boat’.}
		\end{tablenotes}
	\end{threeparttable}
\end{adjustbox}
\end{table}

